\subsection{一张图串起来：从文字到可验证的应用}
\label{app:concept-map}

这张图把模型的\textbf{信息流}与\textbf{学习信号}放在一起：
上方沿前向计算得到概率，中间由训练目标形成梯度，
下方展示后训练、推理生成和外部工具各自的位置。按章号与页码可回读对应推导。

\begin{center}
\tagstructbegin{tag=Figure,alttext={从 Token 到 Embedding、注意力、输出、损失、梯度、参数更新，再到推理和 Agent 工具链的概念关系图}}
\begin{tikzpicture}[
  >=Latex,
  node distance=6mm,
  unit/.style={draw=linegray,fill=paper,rounded corners=3pt,minimum width=25mm,minimum height=13mm,align=center,font=\sffamily\small},
  signal/.style={draw=flowcolor,fill=linegray!15,rounded corners=3pt,minimum width=31mm,minimum height=12mm,align=center,font=\sffamily\small},
  app/.style={draw=flowcolor,fill=linegray!25,rounded corners=3pt,minimum width=31mm,minimum height=12mm,align=center,font=\sffamily\small},
]
\node[unit] (tok) at (0,0) {Token/ID};
\node[unit] (emb) at (3.2,0) {Embedding};
\node[unit] (att) at (6.4,0) {Attention};
\node[unit] (ffn) at (9.6,0) {FFN/输出头};
\foreach \a/\b in {tok/emb,emb/att,att/ffn}{\draw[->,thick,flowcolor] (\a)--(\b);}
\node[signal] (loss) at (9.6,-2.15) {概率与交叉熵\\[-1mm]$\displaystyle L=-\ln p_y$};
\node[signal] (grad) at (6.4,-2.15) {反向传播\\[-1mm]$p-\mathrm{onehot}$};
\node[signal] (upd) at (3.2,-2.15) {更新参数\\[-1mm]$\theta\gets\theta-\eta\nabla L$};
\draw[->,thick,flowcolor] (ffn)--(loss);
\draw[->,thick,flowcolor] (loss)--(grad);
\draw[->,thick,flowcolor] (grad)--(upd);
% The arrowhead must finish at the model, not above or through the Embedding box.
\draw[->,thick,flowcolor] (upd.north) --
  node[midway,right,font=\sffamily\tiny,text=flowcolor] {参数回写} (emb.south);
% Post-training is a separate optimization phase, not a step after inference.
\node[signal] (post) at (1.65,-4.45) {后训练\\行为目标};
\node[app] (gen) at (5.0,-4.45) {推理生成\\温度/缓存};
\node[app] (tools) at (8.7,-4.45) {RAG/工具/Agent\\系统边界};
% Post-training can update parameters; inference can enter an external application.
% Neither route claims that tool calls automatically backpropagate.
\draw[->,densely dashed,flowcolor] (post.north) -- ++(0,5mm) -| (upd.south);
\draw[->,thick,flowcolor] (gen.east) -- (tools.west);
\end{tikzpicture}
\captionof{figure}{模型内部的前向计算、损失反传与外部工具链是三个不同层次}
\tagstructend
\end{center}

\noindent
\textbf{顺着图回读：}
\hyperref[intro:20-embeddings]{Embedding} $\rightarrow$
\hyperref[intro:11-softmax-memory]{Attention} $\rightarrow$
\hyperref[intro:22-transformer-architecture]{Transformer Block} $\rightarrow$
\hyperref[intro:22c-language-training]{训练} $\rightarrow$
\hyperref[intro:23-llm-systems]{生成} $\rightarrow$
\hyperref[intro:23d-agent-systems]{应用系统}。


